\section{Funktionale Anforderungen}
Die funktionalen Anforderungen beschreiben alle Funktionen, die das System bereitstellen muss.
Diese ergeben sich direkt aus der Zielsetzung sowie aus den identifizierten Defiziten bestehender Systeme.

Die zentrale Funktion des Systems besteht in der kontinuierlichen und unterbrechungsfreien Darstellung der Uhrzeit in sprachlicher Form.
Im Gegensatz zu klassischen WordClocks erfolgt zusätzlich eine minutengenaue und sekundengenaue Ergänzung, ohne die Hauptanzeige zu überlagern.
Weiterhin werden externe Datenquellen integriert.
Hierzu zählen insbesondere Wetterinformationen sowie die Außentemperatur.
Diese Daten werden in abstrahierter Form dargestellt, um die sprachliche Konsistenz der Anzeige zu erhalten.

Zusätzliche Sensorik ermöglicht die Erfassung von Innenraumparametern.
Die Anzeige erfolgt über separate \gls{led}-Elemente und ergänzt die Hauptfunktion um informationsorientierte Inhalte.
Die Benutzerinteraktion erfolgt über eine \gls{ui}, welche über eine webbasierte Schnittstelle bereitgestellt wird.
Dadurch wird eine flexible Steuerung ohne zusätzliche Hardware ermöglicht.
Alle funktionalen Anforderungen sind in Tabelle~\ref{tab:Funktionale Anforderung} aufgelistet.

\begin{table}[H]
\caption{Funktionale-Anforderungen der WordClock}
\centering
\begin{tabular}{p{0.5cm} p{15cm}}
\textbf{ID} & \textbf{Beschreibung} \\
\hline
F1 & Wortbasierte Anzeige der aktuellen Uhrzeit (Stunde und Minute) \\
F2 & Kontinuierliche Sekundenanzeige ohne Unterbrechung von F1 \\
F3 & Minutengenaue Darstellung durch zusätzliche \gls{led}-Elemente \\
F4 & Anzeige der Außentemperatur in sprachlicher Form basierend auf Schwellenwerten \\
F5 & Darstellung von Wetterzuständen in sprachlicher Form \\
F6 & Anzeige der Innenraumtemperatur \\
F7 & Anzeige der Luftfeuchtigkeit im Innenraum \\
F8 & Symbolische Darstellung der Tageszeit \\
F9 & Symbolische Darstellung der Mondphase \\
F10 & Abruf externer Daten über Internetverbindung \\
F11 & Benutzerinteraktion zur Steuerung von Anzeige, Helligkeit und Farbe \\
F12 & Bereitstellung einer \gls{ui} zur Systemsteuerung und Darstellung des Betriebszustand \\
\end{tabular}
\label{tab:Funktionale Anforderung}
\end{table}